\section{Introduction}
	\label{introduction}
	
	\par Readers primarily interested in implementation details may refer to subsection \ref{sec:finalImplementation}.\newline
	
	\par This tutorial assumes prior familiarity with neural networks and their underlying mechanisms.
	
	\par Standard neural networks (NNs), herein defined as non-spiking, require activation of all neurons during both forward and backward propagations. This means that every unit in the network must perform computations at every time step, leading to significant power consumption \cite{10242251}. In an era where responsible use of energy resources is essential, a critical question arises: How can one design NN models that are energy-efficient enough to be deployed on portable or resource-constrained devices?
	
	\par Spiking Neural Networks (SNNs) offer a compelling answer to this problem. As the third generation of neural networks, SNNs encode information using discrete spike events and update neuron states only when necessary. This event-driven paradigm mimics biological neuronal dynamics and allows for sparse computations, which are particularly advantageous when implemented on neuromorphic hardware. Such systems can achieve high energy efficiency by leveraging the inherent sparsity of spike-based communication.
	
	\par Despite the availability of tools such as SNNTorch for SNN simulation and training, a comprehensive understanding of their theoretical foundations and low-level implementation remains essential — especially in scenarios where existing frameworks are incompatible with specialized hardware or minimal runtime environments.
	
	\par This paper aims to provide a detailed exposition of the theoretical principles of spiking neurons, culminating in a practical implementation guide